This paper makes several contributions to blink detection in fatigue driving research. First, we formalise blink detection in an epoch structured EEG setting and demonstrate the limitations of naive thresholding approaches that do not account for the underlying epoch structure. Second, we propose a three stage blink detection pipeline comprising epoch screening, sample level threshold estimation, and blink region extraction. To our knowledge, this is among the first frameworks to integrate these three stages explicitly to address the challenges associated with epoch structured EEG recordings. Third, we introduce a robust threshold estimator based on the median and median absolute deviation (MAD) and systematically compare its performance with a mean based alternative, Proposed Mean, and conventional baseline methods, including BLINKER concat~\citep{kleifges2017blinker} and MNE annot. Finally, we provide a comprehensive benchmark of the competing methods across 30, 40, and 60~s epoch durations using two naturalistic driving EEG datasets, consisting of one dataset collected in house and one publicly available dataset~\citep{cao2019multichannel,dari2020unsupervised}. To support reproducibility and facilitate further research, we also developed an automated pipeline for extracting ocular indices from EEG recordings and made it freely available as a Python toolbox called \textit{pyblinker} (TODO). The toolbox is intended to facilitate reproducible blink analysis and encourage more consistent reporting of ocular indices in EEG based fatigue research.

